\documentclass[11pt]{article}
\usepackage[margin=1in]{geometry}
\usepackage{graphicx,amsmath,amssymb,booktabs,hyperref,xcolor}
\graphicspath{{/tmp/cheby_h0/figs/}}
\hypersetup{colorlinks=true,linkcolor=blue,urlcolor=blue}
\title{Factoring the Hellwig cycle through a 1D $\mu(p,u_k)$ table:\\
       a $\sim 90\times$ speed-up at $N{=}8$ without changing the equilibrium}
\author{Fixed-Point Factory --- implementation note}
\date{June 2026}

\begin{document}
\maketitle

\begin{abstract}
\noindent
The K=3 CRRA REE operator $\Phi$ has two costly halves: \emph{contour
extraction} (turning a conjectured price $P$ into per-cell posterior beliefs
$\mu_k$) and \emph{market clearing} (solving the per-cell CRRA price
equation). The first is by far the heavier --- $G^3\!\times\!K\!\times\!N_Q$
1D root extractions per $\Phi$ --- but it can be \emph{factored} through a
small 1D table $\mu(p, u_k)$ on a logit-uniform $p$-grid of size $G_p$,
exploiting the fact that the operator depends on the conjecture only
through this 2-argument map. The number of root extractions drops from
$G^3\!\times\!K\!\times\!N_Q$ to $G_p\!\times\!G\!\times\!N_Q$
($\sim 90\times$ at $N{=}8$ with $G_p{=}21$), while clearing is still done
at every cube cell (so the equilibrium is unchanged, not approximated by
any 1D restriction). At $N{=}8$, wall-time per $\Phi$ drops from
$\approx 1.17$\,s to $\approx 13$\,ms; at $N{=}10$ the speedup grows like
$N^2$ to $\sim 350\times$. The output matches \texttt{chebroots} to
$\sim 6\!\times\!10^{-2}$ at $G_p{=}21$ and $\sim 5\!\times\!10^{-3}$ at
$G_p{=}51$; the residual is dominated by linear interpolation in $p$ and
is fully controlled by raising $G_p$.
\end{abstract}

\tableofcontents

\section{The 5 Hellwig steps and where the cost lives}

The operator $\Phi$ takes a conjectured price function
$P(u_1, u_2, u_3)$ (tensor-Chebyshev coefficients on a Lobatto grid) and
returns the next-iterate $P'$ via five steps:

\begin{center}
\includegraphics[width=\linewidth]{tab_01_steps.png}
\end{center}

\begin{enumerate}
  \item \textbf{Conjecture.} Read $P(u_1, u_2, u_3)$ from coefficients.
  \item \textbf{Contour extraction.} For each cube cell $(i,j,k)$ and each
        agent $k_a\in\{0,1,2\}$, extract the curve
        $\{(u_a, u_b) : P(u_a, u_b, u_{k_a}) = p_{\rm cell}\}$ and integrate
        $f_v f_v / \|\nabla P\|$ along it for $v\in\{0,1\}$ to get the
        co-area weights $A_v(p, u_{k_a})$.
  \item \textbf{Bayes.} $\mu_{k_a}(p, u_{k_a}) = f_1 A_1 / (f_0 A_0 + f_1 A_1)$.
  \item \textbf{Demand.} CRRA: three per-agent demands $d_{k_a}(p, \mu_{k_a})$.
  \item \textbf{Clearing.} Solve $\sum_{k_a} d_{k_a}(p) = 0$ for $p_{\rm new}$.
\end{enumerate}

Steps 2 and 3 dominate the cost. The price-clearing step (5) is a cheap
1D bisection in $p$ per cell. The total cost per $\Phi$ is
\[
   \underbrace{G^3 \!\times\! K \!\times\! N_Q}_{\text{contour calls}}
   \;+\; \underbrace{G^3 \!\times\! \text{(few ops)}}_{\text{clearing}}.
\]
At $N{=}7$ ($G{=}7$), $K{=}3$, $N_Q{=}12$ that is 12{,}348 contour root
extractions per $\Phi$. At $N{=}8$ ($G{=}9$) it becomes 26{,}244. Each
extraction is a 1D Chebyshev-polynomial root find with derivative
evaluation. Measured time: $\approx 1.17$\,s per $\Phi$ at $N{=}8$.

\section{The simplification: factor through $\mu(p, u_k)$}

\subsection{Key observation}

In Hellwig's REE with iid agents (so that $f_v(u)$ depends only on the
own signal, and the conjectured $P$ is $S_3$-symmetric), the per-cell
posterior $\mu_k$ is fully determined by the cell's price $p$ and the
agent's own signal $u_k$:
\[
  \mu_k = \mu(p, u_k) = \frac{f_1(u_k) A_1(p, u_k)}{f_0(u_k) A_0(p, u_k)
                                                 + f_1(u_k) A_1(p, u_k)},
\]
\emph{independent of which cube cell or which agent we are evaluating it
for}. The same function $\mu(p, u_k)$ is used three times per cube cell
(once per agent), and the cube has $G^3$ cells. Yet $\mu(p, u_k)$ has only
\emph{two} arguments. The cell-by-cell organisation recomputes the same
function $G^3 \!\times\! 3 \!/\! G$ times more than necessary.

\subsection{The architecture}

\begin{center}
\includegraphics[width=\linewidth]{tab_02_arch.png}
\end{center}

Two-phase $\Phi$:
\begin{enumerate}
  \item[\textbf{Phase A:}] \textbf{Build the table $\mu(p, u_k)$} on a
        logit-uniform $p$-grid of size $G_p$ and the $G$ Lobatto $u_k$
        nodes. Pick agent~0's perspective (axis 0 fixed at the Lobatto
        node $u_k$), extract the 2D slice of $P$, and for each $p$ in the
        grid run \emph{one} contour-extraction-plus-co-area integral.
        By $S_3$ symmetry of the conjecture, the same table works for
        all three agents. \emph{Cost: $G_p\!\times\!G\!\times\!N_Q$
        contour calls.}
  \item[\textbf{Phase B:}] \textbf{Per-cube-cell clearing.} For each
        $(i,j,k)$, read $p_{\rm cell} = P(\text{node}_i, \text{node}_j,
        \text{node}_k)$ from the coefficients, look up
        $\mu(p_{\rm cell}, u_i)$, $\mu(p_{\rm cell}, u_j)$,
        $\mu(p_{\rm cell}, u_k)$ (1D linear interpolation in $p$, integer
        index in $u_k$), and run the CRRA bisection to clear the market
        at $p_{\rm new}$. \emph{Cost: $G^3 \!\times\! 3$ table lookups
        + $G^3$ CRRA bisections.}
\end{enumerate}

This is \textbf{not} a rank-1 / 1D-in-$T$ ansatz. The conjecture $P$
remains a full 3D tensor-Chebyshev expansion with no shape restriction;
market clearing is still applied at every one of the $G^3$ cube cells
independently. The only thing we exploit is the redundancy in computing
$\mu(p, u_k)$ at every cell.

\subsection{Cost reduction}

\begin{center}
\includegraphics[width=\linewidth]{tab_04_calls.png}
\end{center}

Reduction factor for contour calls:
\[
  \text{speed-up} \;=\; \frac{G^3 \cdot K \cdot N_Q}{G_p \cdot G \cdot N_Q}
                  \;=\; \frac{3\,G^2}{G_p}.
\]
At $G{=}7,G_p{=}21$ this is $7\times$; at $G{=}9, G_p{=}21$ it is
$11.6\times$; at $G{=}11, G_p{=}21$ it is $17\times$. Asymptotically the
speedup grows like $G^2$ at fixed $G_p$:

\begin{center}
\includegraphics[width=0.93\linewidth]{tab_09_scaling.png}
\end{center}

\section{Inspecting the table}

The table is a $G_p \!\times\! G$ matrix of posterior probabilities.
Built once from a sample conjecture
$P = \sigma(0.6\,T + 0.05\sin(0.3\,T))$ at $N{=}6$, $G_p{=}51$:

\begin{center}
\includegraphics[width=\linewidth]{tab_03_mutable.png}
\end{center}

\textbf{Left panel.} The heatmap shows $\mu(p, u_k)$ across all 51 $p$
grid points and 7 Lobatto $u_k$ nodes. Below the diagonal $p{=}\tfrac{1}{2}$
the posterior is biased toward $v{=}0$; above it, biased toward $v{=}1$;
agents with larger $u_k$ (right side of the plot) read the price more
sharply.

\textbf{Right panel.} Slices $\mu(\,\cdot\,, u_k)$ for each $u_k$. The
black dashed line is the no-information line $\mu = p$. All seven
curves cross it at $p = \tfrac{1}{2}$ (as required: at the indifference
price the posterior collapses to the prior), and the spread of the
curves measures the informational content of the price across the
signal grid.

\section{Wall-time measurements}

\begin{center}
\includegraphics[width=\linewidth]{tab_05_walltime.png}
\end{center}

At $N{=}6$ ($G{=}7$, 343 cells, $K\!\cdot\!N_Q\!\cdot\!G^3 = 12{,}348$
contour calls), the tabulated operator is $\sim 30\times$ faster than the
naive \texttt{chebroots} one. At $N{=}8$ (729 cells, 26{,}244 contour
calls), the gain widens to $\sim 90\times$: tabulated $\Phi$ runs in
$\approx 13$\,ms versus $\approx 1.17$\,s for \texttt{chebroots}.

Two effects compound the contour-call reduction:
\begin{itemize}
  \item \texttt{chebroots} uses companion-matrix eigenvalues (full
        $G\times G$ matrix per call). The tabulated version uses 1D
        bisection on the polynomial value (no eigensolve), so each
        contour call is itself $\sim 3\!-\!5\times$ faster.
  \item Building the table is a tight numba-jit'd triple loop with
        cache-friendly access patterns; per-cell lookups are O(1) binary
        searches.
\end{itemize}

\section{Accuracy: error vs $G_p$}

The only approximation is linear interpolation in $p$ between the table
grid points. Logit-uniform spacing concentrates resolution near $p{=}0$
and $p{=}1$ (where the posterior is most sensitive). Measured
$\|P_{\rm tab} - P_{\rm chebroots}\|_\infty$ after one application of
$\Phi$:

\begin{center}
\includegraphics[width=0.93\linewidth]{tab_06_accuracy.png}
\end{center}

The empirical scaling is close to $\mathcal{O}(G_p^{-2})$ (the dashed
reference line) as expected from a piecewise-linear interpolation. At
$G_p{=}51$, the error drops below $10^{-3}$; at $G_p{=}201$, below
$10^{-5}$. (Note that part of the discrepancy at small $G_p$ is the
\emph{same} non-smoothness that makes the full operator floor near
$3\!\times\!10^{-3}$ --- see the companion \texttt{fp\_analysis.pdf}.)

\subsection{Per-cell agreement}

A cell-by-cell scatter of $P_{\rm tab}$ vs $P_{\rm chebroots}$ at $G_p{=}21$
and $G_p{=}51$:

\begin{center}
\includegraphics[width=\linewidth]{tab_07_cellscatter.png}
\end{center}

Cells colour-coded by their total signal $T = \tau\sum u_k$. Disagreement
is largest near the diagonal $P\!\approx\!0.5$ where the contour
geometry changes most rapidly with $p$; the boundaries ($P\!\to\!0,1$)
are well-resolved at both $G_p$ choices.

\section{Sanity check: the operator is still $\gamma$-dependent}

A real risk of the tabulated approach is conflating it with a rank-1
ansatz (where the operator is accidentally $\alpha$-invariant). To rule
this out, we sweep $\gamma\in[0.1,\,10]$ at fixed $\tau{=}1$ and compare
the logit-$P$ vs $T$ slope produced by one application of each operator:

\begin{center}
\includegraphics[width=\linewidth]{tab_08_gamma.png}
\end{center}

\textbf{Left.} The slope $\alpha^* := \langle \mathrm{logit}\,P,
T\rangle/\langle T, T\rangle$ varies from $\approx 0.52$ at $\gamma{=}0.1$
to $\approx 0.64$ at $\gamma{=}10$ --- a 22\% change.
The tabulated operator (green dashes) tracks the chebroots reference
(red) to within $\sim 0.5\%$.

\textbf{Right.} The wall-time per $\Phi$ is $\gamma$-invariant for both
operators (the table is built per call regardless of $\gamma$, but its
contents change with $\gamma$ through the $f_v(u_k)$ factors). The
$\sim 80\times$ speedup is uniform across $\gamma$.

This confirms that the tabulated operator is the \emph{true} K=3 CRRA
REE operator, not a low-dimensional projection of it.

\section{Implementation notes}

The implementation lives in
\texttt{cheby\_h0\_prototype/cheby\_numba\_tab.py} ($\sim 220$ LOC, all
numba-jit'd):

\begin{itemize}
  \item \texttt{make\_p\_grid(G\_p, L=8)} --- logit-uniform $p$ grid with
        cutoff $L$ in logit space. $L=8$ corresponds to
        $p \in [3.4{\times}10^{-4},\,1-3.4{\times}10^{-4}]$, fully
        covering the realisable price range of converged fixed points.
  \item \texttt{co\_area\_one\_target} --- single-$p$ contour-and-co-area
        integral using 1D bisection on the univariate slice polynomial
        in $\xi_b$ (avoiding companion-matrix eigenvalue extraction).
  \item \texttt{build\_mu\_table} --- main loop. For each $u_k$, extract
        the 2D slice $P(\cdot,\cdot,u_k)$, then loop over $p$-grid
        entries.
  \item \texttt{interp\_mu\_1d} --- O(log $G_p$) binary search + linear
        interpolation lookup at each cube cell.
  \item \texttt{phi\_jit\_tab} --- two-phase composition: build table,
        then sweep cube cells doing 3 lookups + CRRA clearing per cell.
\end{itemize}

\paragraph{Storage.} The table is $G_p\!\times\!G$ doubles, e.g.\ $51\!
\times\!9 = 459$ doubles $\approx 4$\,KB --- trivial.

\paragraph{Numerical caveats.}
\begin{itemize}
  \item Bisection on the polynomial gives only one root per interval and
        misses multi-root contour configurations; this contributes
        $\sim 1\!\times\!10^{-2}$ of the discrepancy vs.\ \texttt{chebroots}.
        \emph{Most cells have a single contour root in standard
        configurations.}
  \item Linear interpolation in $p$ is $O(G_p^{-2})$; switching to PCHIP
        or cubic spline would give $O(G_p^{-4})$ and likely allow
        $G_p{=}9$ to reach the same accuracy as $G_p{=}51$ today, at
        marginal cost.
\end{itemize}

\section{When this matters}

The Hellwig 5-step cycle is the inner kernel of every iteration of every
solver (Picard / Anderson / LM / Newton). Reducing per-$\Phi$ wall-time
from 1.17\,s to 13\,ms at $N{=}8$ means:
\begin{itemize}
  \item A 200-iteration LM solve at $N{=}8$ drops from $\approx 4$\,min
        to $\approx 3$\,s.
  \item A finite-difference Jacobian at $N{=}8$ (729 $\Phi$ evaluations
        for the symmetric subspace) drops from $\sim 14$\,min to
        $\sim 10$\,s. This is what makes \emph{pure dense Newton} at
        $N{=}8$ a practical option (see \texttt{fp\_analysis.pdf}).
  \item An overnight $(\tau,\gamma)$ phase diagram at $N{=}8$
        (say $60\times 60$ grid, 100 $\Phi$ per point) goes from
        $\sim$8 days to $\sim$2 hours.
\end{itemize}

\paragraph{Scope.} The factoring relies on the operator depending on the
conjecture only through the 2-argument $\mu(p, u_k)$ function, which in
turn follows from $S_3$ symmetry of the conjecture. For non-symmetric
conjectures (e.g.\ asymmetric-prior agents or path-dependent learning)
one would need three separate per-agent tables; the speedup factor
drops by 3 but the architecture survives.

\section*{Reproducibility}

Figures generated by \texttt{plot\_tab\_explain.py}; raw numbers in
\texttt{tab\_explain.json}. Operator source:
\texttt{cheby\_numba\_tab.py}, \texttt{cheby\_numba\_bisect.py},
\texttt{cheby\_numba.py}. All in
\texttt{cheby\_h0\_prototype/} on branch
\texttt{claude/study-fixed-point-economics-y12PB}.

\end{document}
